% =========================================================
%  Section : Le théorème de Freedman
% =========================================================

\subsection{Freedman(1982)}


Le théorème de Freedman, publié en 1982, est le pendant topologique du théorème de Donaldson.
Là où Donaldson impose des contraintes drastiques sur les formes d'intersection réalisables
par des variétés \emph{lisses}, Freedman montre que dans la catégorie \emph{topologique} ces
contraintes disparaissent presque entièrement~: toute forme bilinéaire symétrique entière
unimodulaire est réalisable.

\begin{theorem}[Freedman, 1982]\label{thm:freedman}
  Soit $Q$ une forme bilinéaire symétrique entière unimodulaire. Alors il existe une
  $4$-variété topologique compacte orientée fermée simplement connexe $M$ dont la forme
  d'intersection est isomorphe à $Q$. De plus~:
  \begin{enumerate}
    \item si $Q$ est impaire, $M$ est unique à homéomorphisme près~;
    \item si $Q$ est paire, il existe exactement deux telles variétés à homéomorphisme près,
          qui se distinguent par leur invariant de Kirby-Siebenmann
          $\mathrm{ks}(M) \in \mathbb{Z}/2\mathbb{Z}$.
  \end{enumerate}
\end{theorem}

C'est un résultat de classification complète~: les $4$-variétés topologiques compactes
orientées fermées simplement connexes sont entièrement déterminées, à homéomorphisme près,
par leur forme d'intersection et leur invariant de Kirby-Siebenmann.


\subsubsection*{Conséquences et portée}

Le théorème de Freedman a des conséquences immédiates sur la classification des $4$-variétés
topologiques surtout, combiné au théorème de Donaldson, il produit la conséquence suivante qui est au
c\oe ur de ce mémoire~:

\begin{corollary}\label{cor:freedman_donaldson}
  Il existe des $4$-variétés topologiques compactes orientées fermées simplement connexes qui
  n'admettent aucune structure différentielle.
\end{corollary}

\begin{proof}
  Par le théorème de Freedman, il existe une $4$-variété topologique $M_{E_8}$ de forme
  d'intersection $Q_{E_8}$. Par le théorème de Donaldson, toute $4$-variété lisse simplement
  connexe à forme définie positive a une forme diagonalisable sur $\mathbb{Z}$. Comme
  $Q_{E_8}$ est définie positive et non diagonalisable (voir Section~\ref{prop:e8_non_diag}),
  $M_{E_8}$ n'admet aucune structure lisse.
\end{proof}

\begin{remark}
  La portée du théorème de Freedman va bien au-delà de ce corollaire. Il a notamment permis
  de démontrer la conjecture de Poincaré topologique en dimension $4$~: toute $4$-variété
  topologique simplement connexe ayant le type d'homotopie de $S^4$ est homéomorphe à $S^4$.
  C'est pour ce résultat, ainsi que pour le théorème de classification, que Freedman a reçu
  la médaille Fields en 1986.
\end{remark}